\section{Parallel layouts}
\label{sec:layouts}

\begin{question}
  \textbf{Which parallel layout should each workload use on this partition?} For Qwen3.8-27B: which layout,
  at which batch per GPU, and which lever (batch, gradient-reduction dtype, pipeline parallelism,
  activation checkpointing) buys the most? How much of the step is spent in exposed communication? For
  GLM-5: should expert parallelism stay inside a node or cross nodes?
\end{question}

\paragraph{Qwen3.8-27B (A-Q).} The factors and levels are in \cref{tab:aq}. The full cross of layouts and
tokens per GPU runs at 13 nodes; the 3 best layouts also run at 2 and 52 nodes, and the layouts without
pipeline parallelism at 1 node. About 45 runs, each repeated twice, about 10 minutes each \est{}.

\begin{table}[h]
  \caption{Factors of the Qwen layout experiment (A-Q).}
  \label{tab:aq}
  \centering\small
  \begin{adjustbox}{max width=\linewidth}
  \begin{tabular}{ll}
    \toprule
    Factor & Levels \\
    \midrule
    Layout & \texttt{fsdp} / \texttt{hsdp8} / \texttt{hsdp16} / \texttt{tp2-fsdp} / \texttt{tp8-fsdp} /
             \texttt{pp2-fsdp} / \texttt{pp4-fsdp} \\
    Tokens per GPU & 12k ($4096\times3$) / 24k ($4096\times6$) / 48k ($4096\times12$) \\
    Gradient reduction & FP32 / BF16 (on the 2 best layouts) \\
    Activation checkpointing & selective / full / none (1 and 13 nodes; none only where it fits) \\
    Nodes & 13 (full cross); 2 and 52 for the 3 best; 1 for the layouts without PP \\
    \bottomrule
  \end{tabular}
  \end{adjustbox}
\end{table}

\begin{figure}[h]
  \centering\fig{qwen_layouts}
  \caption{Qwen3.8-27B throughput (tokens/s per GPU, or MFU) by layout (rows) and tokens per GPU
    (columns), at 2, 13 and 52 nodes (panels). \syntheticnote{qwen_layouts}}
  \label{fig:qwen_layouts}
\end{figure}

\placeholder{best layout and batch at each scale; whether keeping the shards in a node (\texttt{hsdp8},
\texttt{tp8-fsdp}) wins as expected.}

\begin{figure}[h]
  \centering\fig{qwen_levers}
  \caption{Throughput change of each lever (batch, gradient-reduction dtype, pipeline parallelism,
    activation checkpointing) relative to \texttt{fsdp} at 12k tokens per GPU, at 13 nodes.
    \syntheticnote{qwen_levers}}
  \label{fig:qwen_levers}
\end{figure}

\placeholder{ranking of the levers; comparison with the network arithmetic of \cref{sec:method}.}

\paragraph{Where the step time goes.} \Cref{fig:step_decomposition} splits the step time into compute,
taken from 1-node runs at the same batch per GPU, and exposed communication, the remainder of the
measured step time. With 8 GPUs on NVLink, the 1-node runs train the full model, so the compute part is
measured, not scaled from a truncated model.

\begin{figure}[h]
  \centering\fig[0.85\linewidth]{step_decomposition}
  \caption{Step time of Qwen3.8-27B split into compute (1-node step time) and exposed communication, by
    layout and number of nodes. \syntheticnote{step_decomposition}}
  \label{fig:step_decomposition}
\end{figure}

\placeholder{share of exposed communication by layout at 13 and 52 nodes.}

\paragraph{GLM-5 (A-G).} The global batch is 64 sequences per model copy. The layouts differ in where
the expert groups lie (\cref{tab:ag}). Recomputation (full vs.\ selective) and micro-batch size (1 vs.\
2) are varied on the best layout. Each runs at 13 nodes (1 copy) and 52 nodes (4 copies). About 14 runs,
each repeated twice, about 25 minutes each \est{}, including about 15 minutes of startup at scale \est{}.

\begin{table}[h]
  \caption{GLM-5 layouts (A-G).}
  \label{tab:ag}
  \centering\small
  \begin{tabular}{llll}
    \toprule
    Layout & Expert group & Nodes per copy & Copies on 52 nodes \\
    \midrule
    \texttt{pp13-ep8}  & in a node (NVLink) & 13 & 4 \\
    \texttt{pp13-ep16} & across 2 nodes & 13 & 4 \\
    \texttt{pp13-ep32} & across 4 nodes & 13 & 4 \\
    \texttt{pp26-ep8}  & in a node & 26 & 2 \\
    \texttt{pp6-ep16} (13 layers per stage) & across 2 nodes & 12 & 4 (48 nodes), if it fits \\
    \bottomrule
  \end{tabular}
\end{table}

\begin{figure}[h]
  \centering\fig{glm_layouts}
  \caption{GLM-5 throughput (tokens/s per GPU) and all-to-all time by layout, at 13 and 52 nodes.
    \syntheticnote{glm_layouts}}
  \label{fig:glm_layouts}
\end{figure}

\placeholder{cost of expert groups crossing nodes; best layout, recomputation and micro-batch.}
